\section{Functorial translations between layers}\label{sec:functorial}

The APEIRON architecture comprises seventeen layers organised into four clusters.
A key claim of the framework is that semantic translations between layers are naturally captured by \emph{functors} between the corresponding categories.
With the Apeiron category $\mathcal{A}$ established for Layer~1, we can now sketch how these functors arise.

For each higher layer $L$, we can define a category $\mathcal{C}_L$ whose objects are the entities of that layer (e.g., relational hypergraphs for Layer~2, functional entities for Layer~3, etc.) and whose morphisms are structure‑preserving maps.
The emergence from Layer~$k$ to Layer~$k+1$ is then formalised as a functor $F_k:\mathcal{C}_k\to\mathcal{C}_{k+1}$.
The functoriality condition --- $F_k(g\circ f)=F_k(g)\circ F_k(f)$ --- ensures that compositional structure is preserved as abstraction increases.

\Cref{fig:layer_functors} illustrates the network of functors connecting the first few layers.
\begin{figure}[htbp]
\centering
\begin{tikzcd}
\mathcal{C}_1 \arrow[r, "F_{1}"] \arrow[dr, bend right, "G_{1}"'] & \mathcal{C}_2 \arrow[r, "F_{2}"] & \mathcal{C}_3 \arrow[r, "F_{3}"] & \dots \\
& \mathcal{C}_4 \arrow[ur, bend right, "G_{3}"'] & &
\end{tikzcd}
\caption{Excerpt of the functor network. $F_k$ are the primary emergence functors; $G_k$ are downward reflection functors (e.g., from functional entities back to relational hypergraphs).}\label{fig:layer_functors}
\end{figure}

Theorem~\ref{th:separated} guarantees that an atomic observable, being separated in $\mathcal{A}$, is mapped by any faithful functor to a separated object in the higher layer, thereby preserving its foundational status throughout the layered hierarchy.

\subsection{Preservation of separatedness}
Because separated objects have no non‑trivial outgoing morphisms, any functor $F_k : \mathcal{C}_k \to \mathcal{C}_{k+1}$ that is faithful on morphisms will preserve separatedness: if $o$ is separated in $\mathcal{C}_k$, then $F_k(o)$ is separated in $\mathcal{C}_{k+1}$.
This ensures that the fundamental irreducibility of observables is maintained from the substrate all the way up to the integrative layers.

\subsection{Toward adjunctions between layers}
In categorical semantics, the most robust form of translation between abstraction levels is an \emph{adjunction} $L \dashv R$, where the left adjoint $L$ freely generates the higher‑layer structure and the right adjoint $R$ extracts the most faithful lower‑layer approximation.
For instance, the emergence of a relational hypergraph (Layer~2) from raw observables (Layer~1) can be seen as a functor $L_{1,2} : \mathcal{C}_1 \to \mathcal{C}_2$ whose right adjoint $R_{2,1}$ would reconstruct the co‑activation probabilities from the hypergraph topology.
Proving the adjunction $L_{1,2} \dashv R_{2,1}$ would formally guarantee that no spurious relations are introduced and that all genuine co‑occurrences are preserved.
We leave the systematic investigation of adjunctions between all seventeen layers for future work.